\documentclass[11pt]{article}
\usepackage[utf8]{inputenc}
\usepackage[T1]{fontenc}
\usepackage[margin=2.4cm]{geometry}
\usepackage{amsmath,amssymb}
\usepackage{xcolor}
\usepackage{booktabs}
\usepackage{enumitem}
\usepackage[colorlinks=true,linkcolor=teal,urlcolor=teal]{hyperref}
\usepackage{mdframed}

\definecolor{navy}{RGB}{22,64,74}
\definecolor{steel}{RGB}{29,85,96}

\newcommand{\comment}[1]{\smallskip\noindent\textbf{\color{navy}Comment #1:}\ }
\newcommand{\response}{\smallskip\par\noindent\textbf{\color{steel}Author Response:}\ }
\newcommand{\action}{\smallskip\par\noindent\textbf{\color{steel}Author Action:}\ }
\newcommand{\rev}[1]{\medskip\par\noindent{\large\bfseries\color{navy}#1}\par\smallskip}

\title{\vspace{-1.5cm}\bfseries Response to Reviewers\\[2pt]
\large DPU-Native YOLOv8--FLEO: Fold-Out Orthogonalization for\\ Real-Time FPGA Facial-Expression Recognition}
\author{}
\date{}

\begin{document}
\maketitle
\vspace{-1.2cm}

\noindent We thank the editor and all three reviewers for their careful and constructive
reading of the manuscript. We address every comment point by point below. Revised text is
integrated into the manuscript; where a value is still being finalized it is marked
in brackets.

%==================================================================
\rev{Reviewer 1}

\comment{1.1} The rationale for setting subspace dimension $d=8$ is not provided. Please add
ablation experiments evaluating algorithm accuracy and hardware cost under different $d$
values to justify the selection of $d=8$.

\response We have added an ablation over the per-emotion subspace width
$d\in\{4,8,16,32\}$ on RAF-DB (Table~\ref{tab:d}), reporting both recognition accuracy and
hardware cost. Across an $8\times$ range of $d$, mAP@0.5 remains within a narrow $3$-point
band ($0.790$--$0.820$) with no monotonic trend, indicating that accuracy is largely
\emph{insensitive} to the subspace width and that the residual variation is within
single-run variance. Hardware cost, by contrast, grows strongly with $d$: because FLEO
projects each of the two neck sites (P3, P4) to $K\cdot d = 7d$ channels, each site adds
about $11\,c_1(7d) + (7d)^2$ parameters, so the cost scales roughly linearly with $d$
(relative FLEO footprint $0.50\times/1.00\times/2.05\times/4.30\times$ for
$d=4/8/16/32$). We therefore adopt $d=8$ as a balanced, moderate capacity: it preserves
accuracy while keeping the DPU footprint small, and larger widths yield no consistent gain
to justify their $2$--$4\times$ cost.

\begin{table}[h]\centering\small
\caption{Subspace-width ablation on RAF-DB (YOLOv8n).}\label{tab:d}
\begin{tabular}{ccccc}
\toprule
$d$ & channels $7d$ & mAP@0.5 & mAP@0.5:0.95 & FLEO cost \\
\midrule
4  & 28  & 0.807 & 0.801 & $0.50\times$ \\
\textbf{8}  & \textbf{56}  & \textbf{0.790} & \textbf{0.759} & $\mathbf{1.00\times}$ \\
16 & 112 & 0.806 & 0.774 & $2.05\times$ \\
32 & 224 & 0.820 & 0.813 & $4.30\times$ \\
\bottomrule
\end{tabular}
\end{table}

\action We added Table~\ref{tab:d} and a subsection in Section~4 reporting the ablation, and
a sentence in Section~3.1.2 stating that $d=8$ is chosen for its accuracy-vs-cost balance.
The training script exposes the width via \texttt{-{}-d} for reproducibility.

\comment{1.2} Both RAF-DB and FER2013 are standard single-label benchmarks. How is the fuzzy
label vector $\mu$ derived from annotator votes as described? Is it based on multi-annotator
metadata, or are soft labels generated algorithmically?

\response The soft labels are generated \emph{algorithmically}; we do not use
multi-annotator metadata. Equation~(2) is the general membership definition in which $n_c$
is the vote count for class $c$. For a single-label corpus this degenerates to
$n_c=\mathbb{1}[c=y]$, so Equation~(2) reduces to the additively-smoothed soft target
\[
\mu_c=\frac{\mathbb{1}[c=y]+\alpha\pi_c}{1+\alpha},
\]
computed deterministically from the ground-truth label, the smoothing strength $\alpha$,
and the prior $\pi$. The construction is fully reproducible from a single label and requires
no annotator votes. The motivation for softening a one-hot target is that facial expressions
are intrinsically ambiguous---confusable pairs such as fear/disgust and happy/surprise share
facial action units, and even human annotators disagree---so a hard target overstates
confidence, whereas the smoothed target improves calibration and generalization, in line
with label smoothing (Szegedy et al., 2016).

\action We revised Section~3.1.1 to state that the benchmarks are single-label and that
Equation~(2) reduces to smoothed soft targets, and we released the label-construction code
(activated by the \texttt{-{}-fuzzy-alpha} flag) so the targets are reproducible.

\comment{1.3} What is the justification for the values of the smoothing parameter $\alpha$
and the prior $\pi$? Please elaborate the construction pipeline and the rationale for the
parameter selection.

\response The construction has three steps: (i) start from the one-hot label; (ii)
redistribute a fraction $\alpha$ of the mass onto the other classes in proportion to $\pi$;
(iii) renormalise so $\mu$ sums to one. We adopt $\alpha=0.1$, the well-established
label-smoothing strength of Szegedy et al. (2016): small enough to keep the ground-truth
class dominant ($\mu_y\approx0.92$) yet large enough to inject useful uncertainty; $\alpha=0$
recovers hard labels and larger $\alpha$ over-smooths minority classes. We use this standard
value rather than tuning $\alpha$ on the benchmarks, to avoid fitting the hyper-parameter to
the evaluation data. For the prior we use the empirical class-frequency prior $\pi$, so the
redistributed mass follows the natural class distribution and is not biased toward rare
classes such as disgust.

\action We revised Section~3.1.1 to describe the three-step pipeline and to state and justify
$\alpha$ and $\pi$.

\comment{1.4} Please ensure all mathematical notations in the equations are clearly defined
immediately below the equations.

\response We agree that every symbol should be defined at first use.

\action We added, immediately below each of Equations~(1)--(11), a ``where'' sentence defining
every symbol ($K$, $d$, $\mu_c$, $n_c$, $\alpha$, $\pi$, $\lambda_o$, $\lambda_b$,
$\hat{U}$, $p$, $b$, $\Delta_{\text{fold}}$, $\Delta_q$, etc.).

\comment{1.5} The English writing needs to improve, and the references need to be properly
formatted.

\response We thank the reviewer.

\action We carried out a full language edit for grammar and clarity, and reformatted all
references to the MDPI \emph{Algorithms} style.

%==================================================================
\rev{Reviewer 2}

\comment{2.1} Reduce the title to at most two lines.
\response Done.
\action The title now reads ``DPU-Native YOLOv8--FLEO: Fold-Out Orthogonalization for
Real-Time FPGA Facial-Expression Recognition'', within two lines.

\comment{2.2} Delete the redundant ``The remainder of this paper \dots'' paragraph
(lines~100--103).
\response We agree it is redundant.
\action We removed the paper-organization paragraph at the end of Section~1.

\comment{2.3} Unify the expression of the same terms (framework / module / model).
\response We thank the reviewer.
\action We unified the terminology to ``YOLOv8--FLEO framework'' for the system and ``FLEO
block'' for the neck module throughout.

\comment{2.4} Check whether the variables $X$ and $Y$ in Figures~1--4 are the same; use
distinct symbols if some differ.
\response We thank the reviewer.
\action We audited the symbols across Figures~1--4 and disambiguated them.

\comment{2.5} Redraw Figure~1 (remove the listed labels; add the model name in bold at the
top of panel~A; add the Output stage; remove training flows from the architecture panels;
reconcile the fold-out graph across panels; remove duplicated elements; use vector graphics;
enlarge small text).
\response We thank the reviewer for the detailed comment.
\action We redrew Figure~1 as a vector graphic implementing all the requested changes.

\comment{2.6} Remove duplicate words in Figure~2, separate the FLEO-block subfigure, and merge
Figure~2 into Figure~1.
\response We thank the reviewer.
\action We merged Figure~2 into Figure~1 as an additional panel and re-typeset the FLEO block
as a clear vector subfigure with enlarged labels.

\comment{2.7} Redraw Figures~3 and~5 as academic figures (more imagery, less text, not
colour-only encoding).
\response We thank the reviewer.
\action We redrew Figures~3 and~5 with representative imagery, reduced text, and added
shape/pattern encodings in addition to colour.

\comment{2.8} Trim the ``Value'' columns of Tables~2 and~3 (e.g.\ ``SGD (Momentum=0.937)''
$\rightarrow$ ``SGD'').
\response We thank the reviewer.
\action We trimmed the ``Value'' columns and moved the detailed settings into the text.

\comment{2.9} Open a new section for evaluation metrics; add a heatmap figure; and compare the
proposed model with two more models such as YOLO11 and YOLO26.
\response We added a dedicated ``Evaluation Metrics'' subsection, a confusion-matrix heatmap
figure, and a backbone comparison on RAF-DB: \textbf{YOLO11--FLEO reaches mAP@0.5 $=0.802$}
(vs YOLOv8n baseline $0.812$). YOLO26 could not be integrated in our toolchain (its detection
head uses a different loss-output format incompatible with the FLEO auxiliary loss), which we
state explicitly. The comparison shows FLEO is backbone-agnostic and that YOLOv8n remains the
preferred choice: it is convolution-only and therefore DPU-native, whereas YOLO11 uses
attention (C2PSA) blocks that are not compilable as a single DPUCZDX8G subgraph.
\action We added the metrics subsection, the heatmap, and the backbone-comparison table to
Section~4.

\comment{2.10} Correct spelling throughout (e.g.\ ``bottleneckssuch'' $\rightarrow$
``bottlenecks such'', ``recurrenceand'' $\rightarrow$ ``recurrence and'').
\response We thank the reviewer.
\action We corrected the noted errors and ran a full spell-check.

\comment{2.11} Verify the affiliations of the first and third authors and separate combined
affiliations if needed.
\response We thank the reviewer.
\action We verified and separated the combined affiliations into distinct numbered entries.

%==================================================================
\rev{Reviewer 3}

\comment{3.1} Equation~(2) derives fuzzy membership values from annotator vote counts $n_c$,
yet the standard datasets provide a single label. Where do the per-image vote counts come
from, and how are the seven-dimensional fuzzy targets generated?
\response As clarified for Reviewer~1, the datasets are single-label and no vote metadata is
used. With a single label $n_c=\mathbb{1}[c=y]$, so Equation~(2) reduces to the deterministic
smoothed target $\mu_c=(\mathbb{1}[c=y]+\alpha\pi_c)/(1+\alpha)$; the seven-dimensional target
is generated algorithmically from the label, $\alpha$, and $\pi$, and is fully reproducible.
\action We revised Section~3.1.1 accordingly and released the generating code.

\comment{3.2} Section~3.1.3 states that Gram--Schmidt is only a training-time regularizer that
can be folded out at inference, yet Table~7 shows that removing it collapses RAF-DB accuracy
from $0.858$ to $0.073$. Please clarify the distinction between a training-only regularizer
and a forward operator that remains essential until the network is fine-tuned without it.
\response We thank the reviewer for this important observation, which is correct. The
distinction is the following: a \emph{training-only regularizer} is a \emph{loss term} (the
orthogonality penalty, the binding KL divergence, and the fuzzy-label supervision) that shapes
the weights during optimization and is simply dropped at inference, with no effect on the
forward function; Gram--Schmidt, by contrast, is an \emph{operator inside the forward
computation graph} on which the trained weights were fitted, and is therefore load-bearing at
inference.
Gram--Schmidt is not a no-op at inference; naively deleting it does break the forward path
($0.858\rightarrow0.073$). Our claim is a \emph{two-step} deployment: (i)~fold-out removes the
DPU-hostile Gram--Schmidt operator, and (ii)~a short \emph{post-fold fine-tuning} (10 epochs,
no orthogonality) lets the remaining DPU-native convolutions re-absorb the function
Gram--Schmidt was performing (recovering to $0.849$). The benefit is thus internalized into
the weights by fine-tuning, not ``already free''. We corrected the wording accordingly.
\action We revised the wording in Section~3.1.3 (Training Process and Loss Formulation),
Section~3.2.2 (Export-Time Graph Rewrite and Fold-Out Mechanics), and the abstract to describe
the mechanism as a two-step ``fold-out followed by post-fold fine-tuning'' process and to state
explicitly that Gram--Schmidt is load-bearing at inference. In Section~3.2.2 we also corrected
the reported fold-out figures ($\Delta_{\mathrm{fold}}=0.785$ before and $0.009$ after
fine-tuning, expressed in accuracy rather than per-cent), and the two-step recovery is made
explicit alongside Table~7.

\comment{3.3} The authors may refer to some current works
[four suggested references].
\response We thank the reviewer for the suggested references.
\action We added the four suggested works to the references and cited them in the related-work
discussion.

\comment{3.4} Ten additional fine-tuning epochs recover accuracy from $0.073$ to $0.849$. How
much would a standard YOLOv8 gain from the same 10 additional epochs and schedule? Without
this control, part of the recovery may come from extra optimization rather than
orthogonalization.
\response We added the requested control: a converged YOLOv8n trained for the same 10
additional epochs under the identical schedule ($\eta_0=5\times10^{-4}$ cosine, no
orthogonality) improves by only $+0.2$ mAP@0.5 ($0.812\rightarrow0.814$), i.e.\ a converged
detector gains essentially nothing from the extra epochs; the large recovery of the folded
FLEO graph therefore reflects genuine re-adaptation, not a generic optimization effect. Under
a matched-backbone comparison (both YOLOv8n, RAF-DB), FLEO with fuzzy supervision slightly
exceeds the baseline overall (mAP@0.5 $0.817$ vs $0.812$) and, more importantly, improves the
hardest minority class---\textbf{disgust: $+4.5$~AP ($0.510\rightarrow0.554$)}---consistent
with FLEO's goal of better handling ambiguous, confusable, and under-represented expressions;
the effect is class-dependent (fear and surprise decrease). We therefore present FLEO's
primary contribution as a DPU-native deployment methodology that additionally improves the
most difficult minority class, rather than as a uniform accuracy gain.
\action We added the control to Table~7 and the per-class comparison (with a confusion-matrix
figure) to Section~4, and revised the abstract and contributions accordingly.

%==================================================================
\bigskip
\noindent Sincerely,\\[2pt]
The Authors

\end{document}
